The clique complex \(\Cl(H)\) contains the cliques of \(H\), including
the empty face. We use \defi{augmented reduced homology over \(\F\)}. For a complex \(K\), the chain group
\(\Ct_t(K)\) has the faces of cardinality \(t+1\) as a basis when
\(t\geq0\), and
\(\Ct_{-1}(K)=\F\), generated by the empty face. Lower chain groups
are zero. The boundary of a face is the sum of its codimension-one
faces; in particular the boundary of a vertex is the empty face.
Each term occurs twice in \(\partial^2\), so this is a chain complex.

We write \(\Ht_t(K)=\ker\partial_t/\operatorname{im}\partial_{t+1}\)
and call \(K\) \defi{acyclic} if every \(\Ht_t(K)\) vanishes.
In particular,
\begin{equation}\label{hom:eq:empty}
 \Ht_{-1}(\{\varnothing\})=\F,\qquad
 \Ht_t(\{\varnothing\})=0\quad(t\ne-1).
\end{equation}
A complex with a vertex has zero homology in degree \(-1\).
Its \defi{reduced Betti number} in degree \(t\) is
\(\widetilde\beta_t(K)=\dim_{\F}\Ht_t(K)\).
Acyclicity is used as a homological condition, not as a synonym for
contractibility or for membership in any \(\D_j\).
When a homological term is applied to a graph, it refers to its clique
complex; in particular, \emph{acyclic} here does not mean that the graph is a forest.

A \defi{simplicial cone} with apex \(w\) is a complex
\(w*J\) whose faces are exactly the faces \(\sigma\) of \(J\)
and the sets \(\sigma\cup\{w\}\), where \(w\notin V(J)\).
In particular, if a graph \(H\) is a cone with apex \(w\), then
\(\Cl(H)=w*\Cl(H-w)\).

\begin{lemma}\label{hom:lem:cone}
A simplicial cone is acyclic.
\end{lemma}
\begin{proof}
For an apex \(w\), let \(h\) adjoin \(w\) to a face not containing it
and send a face containing \(w\) to zero. Direct expansion gives
\(\partial h+h\partial=\mathrm{id}\), including on the empty face.
Every cycle \(z\) is therefore the boundary of \(h(z)\).
\end{proof}

We will compare complexes by expressing them as unions of acyclic pieces.

\begin{lemma}\label{hom:lem:union}
If \(K=U\cup V\) and \(U,V\) are acyclic, then for every integer \(t\),
\[
 \Ht_t(K)\cong\Ht_{t-1}(U\cap V).
\]
The intersection may be \(\{\varnothing\}\).
\end{lemma}
\begin{proof}
For \(t\geq0\), write a cycle \(z=u+v\), with chains in \(U,V\).
Then \(\partial u=\partial v=b\) is a cycle in \(U\cap V\).
Changing the decomposition adds the same chain of the intersection
to \(u\) and \(v\), so it changes \(b\) by a boundary. If \(z\) is
changed by a boundary, decompose a filling into chains of \(U,V\);
the resulting class of \(b\) is unchanged. This defines
\([z]\mapsto[b]\).

Every cycle \(b\) of the intersection has fillings \(u\) in \(U\)
and \(v\) in \(V\), since both complexes are acyclic; then \(u+v\)
proves surjectivity. If \(b=\partial a\) in the intersection, the
chains \(u+a,v+a\) are cycles in \(U,V\) and hence boundaries.
Their sum is \(z\), proving injectivity.
For \(t=-1\), the left side is zero since the acyclic complexes
\(U,V\) have vertices; the right side is in degree \(-2\).
In lower degrees both sides vanish.
\end{proof}

The folds used in the graph reductions preserve these homology groups.

\begin{lemma}\label{hom:lem:fold}
A fold \(v\to w\) induces isomorphisms
\(\Ht_t(\Cl(H-v))\cong\Ht_t(\Cl(H))\) for every integer \(t\).
\end{lemma}
\begin{proof}
Pair every face \(\sigma\) containing \(v\) but not \(w\) with
\(\tau=\sigma\cup\{w\}\). The domination condition guarantees that
\(\tau\) is a face. Remove these pairs in decreasing order of
\(|\sigma|\). When a pair is removed, \(\tau\) is the only proper
coface of \(\sigma\) still present: every larger coface was removed
in an earlier pair. The remaining faces form a subcomplex.

To see that removing one such pair preserves augmented homology,
let \(K'\) be the resulting subcomplex of the current complex \(K\)
and put \(r=\dim\sigma\). A cycle of degree \(r+1\) has coefficient
zero at \(\tau\), since its boundary has coefficient zero at
\(\sigma\). For a cycle of degree \(r\), adding its coefficient at
\(\sigma\) times \(\partial\tau\) removes \(\sigma\) without changing
its homology class. In other degrees no basis element is removed.
Thus every class has a representative in \(K'\).

Conversely, suppose a cycle \(z'\) in \(K'\) bounds a chain \(y\)
in \(K\). If \(y\) has degree \(r+1\), its coefficient at \(\tau\)
is zero, because the coefficient of \(\sigma\) in \(z'\) is zero.
If \(y\) has degree \(r\), add its coefficient at \(\sigma\) times
\(\partial\tau\); this removes \(\sigma\) and preserves its
boundary. The resulting filling lies in \(K'\). This proves
injectivity as well, including in degree \(-1\).
After all pairs are removed the complex is exactly \(\Cl(H-v)\).
\end{proof}

The recursive neighbourhood condition therefore gives a homological obstruction at every level.

\begin{lemma}\label{hom:lem:acyclic}
If \(H\in\D_j\), \(j\geq-1\), then \(\Cl(H)\) is acyclic.
\end{lemma}
\begin{proof}
Cones are covered by Lemma~\ref{hom:lem:cone}. Induct on \(j\) and,
at a fixed level, on the number of vertices. The one-vertex graph
is a cone. For a valid first deletion \(v\), the two induction
hypotheses give acyclicity of \(\Cl(H-v)\) and \(\Cl(L_H(v))\).
The exact decomposition
\[
 \Cl(H)=\Cl(H-v)\cup\bigl(v*\Cl(L_H(v))\bigr)
\]
has acyclic pieces and intersection \(\Cl(L_H(v))\).
Lemma~\ref{hom:lem:union} proves the assertion.
\end{proof}

